\documentclass[11pt, oneside]{amsart}   	% use "amsart" instead of "article" for AMSLaTeX format
\usepackage{geometry}                		% See geometry.pdf to learn the layout options. There are lots.
\geometry{letterpaper}                   		% ... or a4paper or a5paper or ...
%\geometry{landscape}                		% Activate for rotated page geometry
%\usepackage[parfill]{parskip}    		% Activate to begin paragraphs with an empty line\usepackage{graphicx}				% Use pdf, png, jpg, or eps§ with pdflatex; use eps in DVI mode
								% TeX will automatically convert eps --> pdf in pdflatex

%\usepackage{{Gratzer-Color-Scheme}	% Active to color theorems red and lemmas blue
\usepackage{amsmath}
\usepackage{centernot}
\usepackage{amsthm}
\usepackage{mathtools}
\usepackage{xcolor}
\usepackage{cleveref}
\usepackage{amssymb}
\usepackage{csquotes}
\usepackage{changepage}

\theoremstyle{definition}
\newtheorem{theorem}{Theorem}
\newtheorem{problem}{Problem}
\newtheorem{lemma}{Lemma}

\newcommand{\inc}[1]{#1_{{+}{+}}}
\newcommand{\N}{\mathbb{N}}
\newcommand{\Z}{\mathbb{Z}}
\newcommand{\Q}{\mathbb{Q}}
\newcommand{\R}{\mathbb{R}}
\newcommand{\sube}{\subseteq}
\newcommand{\supe}{\supseteq}

\newcommand{\contrad}{\Rightarrow\!\Leftarrow}

\renewcommand{\epsilon}{\varepsilon}
\renewcommand{\qedsymbol}{\small\sl Q.E.D.}

\setlength{\parindent}{0em}
\setlength{\parskip}{1em}

\title{Fourier Series}
\author{Joseph Geis}
\author{Wilbert Kim}
\date{2026 August 10}

\begin{document}
\maketitle

\section{Introduction}

\include{introduction.tex}

\section{Exercises}

\begin{problem}
    The problem of Fourier Series is very trivial. This is left as an exercise to the reader. As a reminder,
    \[1+2=3 \in \N\]
    \begin{equation}
        \Omega = \left\{ \int_{-\infty}^{\infty} \frac{\partial^2}{\partial x^2} \exp\left( -\frac{ \left( \sum_{i=1}^{n} x_i \cdot \prod_{j=1}^{m} y_{ij} \right)^{ \frac{1}{p} } }{ \left[ \frac{\sqrt{x + \sqrt{y}}}{\mathbf{M}_{ \mu \nu } \cdot \mathbb{A}^{ \dagger } } \right] } \right) dx \right\}^{\frac{-1}{\infty \cdot \aleph_0}}
    \end{equation}
\end{problem}

%!TEX root = project.tex
\begin{problem}
    Using trigonometric identities when necessary, verify the following integrals.
    \begin{enumerate}[(a)]
        \item For all \(n \in \N\)
        \begin{equation}
        	\int_{-\pi}^{\pi} \cos(nx) ~dx = 0
        	\text{ and }
       		\int_{-\pi}^{\pi} \sin(nx) ~dx = 0
        \end{equation}
        
        \begin{proof}
        	Consider that \(D_x \sin(nx) = n\cos(nx)\), thus \(D_x\left[\frac{1}{n}\sin(nx)\right] = \cos(nx)\) and \(\sin(n\pi) = 0\) for all \(n \in \N\). Then it
			follows that
			\begin{align}
			\int_{-\pi}^{\pi} \cos(nx) ~dx & = \left.\frac{1}{n}\sin(nx)\right|_{x=-\pi}^{x=\pi} \\
			& = \frac{1}{n}\left[ \sin(n\pi) - \sin(-n\pi) \right] \\
			& = \frac{1}{n}\left[ \sin(n\pi) + \sin(n\pi) \right] \\
			& = \frac{1}{n} \cdot 0 \\
			& = 0
			\end{align}
			
			Then, consider that \(D_x \cos(nx) = -n\sin(nx)\), thus \(D_x\left[-\frac{1}{n} \cos(nx)\right] = \sin(nx)\) and \(\cos(-t) = \cos(t)\) for all \(t \in \R\). Then it
			follows that
			\begin{align}
			\int_{-\pi}^{\pi} \sin(nx) ~dx & = \left.-\frac{1}{n}\cos(nx)\right|_{x=-\pi}^{x=\pi} \\
			& = -\frac{1}{n}\left[ \cos(n\pi) - \cos(-n\pi) \right] \\
			& = -\frac{1}{n}\left[ \cos(n\pi) - \cos(n\pi) \right] \\
			& = -\frac{1}{n} \cdot 0 \\
			& = 0
			\end{align}
    	\end{proof}
	
        \item For all \(n \in \N\)
        \begin{equation}
        	\int_{-\pi}^{\pi} \cos^2(nx) ~dx = \pi
			\text{ and }
			\int_{-\pi}^{\pi} \sin^2(nx) ~dx = \pi
        \end{equation}
        
        \begin{proof}
        Consider that \(\cos^2(t) = \frac{1}{2}(1 + \cos{2t})\), \(D_x\left[\frac{1}{2n}\sin(2nx)\right] = \cos(2nx)\), and \(\sin(2n\pi) = 0\) for all \(n \in \N\), thus
        \begin{align}
        	\int_{-\pi}^{\pi} \cos^2(nx) ~dx & = \frac{1}{2}\int_{-\pi}^{\pi} (1 + \cos{2nx}) ~dx \\
			& = \frac{1}{2}\left[ x + \frac{1}{2n}\sin(2nx) \right]^\pi_{-\pi} \\
			& = \frac{1}{2}\left[ \pi - (-\pi) \right] \\
			& = \pi
        \end{align}
        
        Then, consider that \(\sin^2(t) = \frac{1}{2}(1 - \cos{2t})\), thus
        \begin{align}
        	\int_{-\pi}^{\pi} \sin^2(nx) ~dx & = \frac{1}{2}\int_{-\pi}^{\pi} (1 - \cos{2nx}) ~dx \\
			& = \frac{1}{2}\left[ x - \frac{1}{2n}\sin(2nx) \right]^\pi_{-\pi} \\
			& = \frac{1}{2}\left[ \pi - (-\pi) \right] \\
			& = \pi
        \end{align}

        \end{proof}
        
        \item For all \(m, n \in \N\)
        \begin{equation} \label{prob2c:first}
        	\int_{-\pi}^{\pi} \cos(mx)\sin(nx) ~dx = 0
        \end{equation}
        For \(m \neq n\)
        \begin{equation} \label{prob2c:second}
        	\int_{-\pi}^{\pi} \cos(mx)\cos(nx) ~dx = 0
			\text{ and }
			\int_{-\pi}^{\pi} \sin(mx)\sin(nx) ~dx = 0
        \end{equation}
        
        \begin{proof} For the integral in \Cref{prob2c:first}, we integrate by parts by choosing
        \begin{equation}
        \begin{array}{c c}
        u = \sin(nx) & v = \frac{1}{m}\sin(mx) \\
        du = n\cos(nx) ~dx & dv = \cos(mx) ~dx
        \end{array}
        \end{equation}
        
        Then,
        \begin{equation}
        	\int_{-\pi}^{\pi} \cos(mx)\sin(nx) ~dx = \left[\frac{1}{m}\sin(nx)\sin(mx)
				- \frac{n}{m} \int \cos(mx)\sin(nx) ~dx \right]_{-\pi}^{\pi}
        \end{equation}
        
        We integrate by parts again by choosing
        \begin{equation}
        \begin{array}{c c}
        u = \cos(nx) & v = -\frac{1}{m}\cos(mx) \\
        du = -n\sin(nx) ~dx & dv = \sin(mx) ~dx
        \end{array}
        \end{equation}
        
        Then,
        \begin{equation}
        \begin{aligned}
        	\int_{-\pi}^{\pi} \cos(mx)\sin(nx) ~dx  = & \left[ \frac{1}{m}\sin(nx)\sin(mx)\right. \\
					& \left.- \frac{n}{m} \left( -\frac{1}{m}\cos(nx)\cos(mx) - \frac{n}{m} \int \cos(mx)\sin(nx) ~dx \right) \right]_{-\pi}^{\pi}
		\end{aligned}
        \end{equation}
        
        If we let \(A = \int_{-\pi}^{\pi} \cos(mx)\sin(nx) ~dx\), then
        \begin{gather}
        A =  \left[ \frac{1}{m}\sin(nx)\sin(mx) + \frac{n}{m^2}\cos(nx)\cos(mx)\right]_{-\pi}^{\pi} + \frac{n^2}{m^2} A \\
        \left(1 - \frac{n^2}{m^2}\right)A = \left[ \frac{1}{m}\sin(nx)\sin(mx) + \frac{n}{m^2}\cos(nx)\cos(mx)\right]_{-\pi}^{\pi}
        \end{gather}
        and since \(\sin(n\pi) = 0\) for all \(n \in \Z\),
        \begin{align}
        \left(1 - \frac{n^2}{m^2}\right)A & = \left. \frac{n}{m^2}\cos(nx)\cos(mx)\right|_{-\pi}^{\pi} \\
        & = \frac{n}{m^2} \left(\cos(n\pi)\cos(m\pi) - \cos(-n\pi)\cos(-m\pi) \right) \\
        & = \frac{n}{m^2} \left(\cos(n\pi)\cos(m\pi) - \cos(n\pi)\cos(m\pi) \right) = 0
		\end{align}
		
		Hence,
		\begin{equation*}
		A = \int_{-\pi}^{\pi} \cos(mx)\sin(nx) ~dx = 0
		\end{equation*}
		
		For the first integral in \Cref{prob2c:second} (assuming \(m \neq n\)) consider that \(\cos\alpha \cos\beta = \frac{1}{2}[\cos(\alpha + \beta) + \cos(\alpha - \beta)]\). Then,
		\begin{align}
			\int_{-\pi}^{\pi} \cos(mx)\cos(nx) \, dx 
			& = \frac{1}{2}\int_{-\pi}^{\pi} \cos((m+n)x) + \cos((m-n)x) \, dx \\
			& = \frac{1}{2}\left[ \frac{1}{m+n}\sin((m+n)x) + \frac{1}{m-n}\sin((m-n)x) \right]_{-\pi}^{\pi} \\
			& = 0
		\end{align}
		
		For the second integral in \Cref{prob2c:second} (assuming \(m \neq n\)), consider that \(\sin\alpha \sin\beta = \frac{1}{2}[\cos(\alpha + \beta) - \cos(\alpha - \beta)]\). Then,
		\begin{align}
			\int_{-\pi}^{\pi} \sin(mx)\sin(nx) \, dx 
			& = \frac{1}{2}\int_{-\pi}^{\pi} \cos((m+n)x) - \cos((m-n)x) \, dx \\
			& = \frac{1}{2}\left[ \frac{1}{m+n}\sin((m+n)x) - \frac{1}{m-n}\sin((m-n)x) \right]_{-\pi}^{\pi} \\
			& = 0
		\end{align}
        \end{proof}
    \end{enumerate}
\end{problem}

\begin{problem}
\label{exercise3}
    Fourier introduced a formulation for representing some function $f(x)$ in terms of sines and cosines given suitable coefficients $(a_n)$ and $(b_n)$: 
    \begin{equation}
    \label{problem3:1}
        f(x) = a_0 + \sum_{n=1}^{\infty} a_n \cos(nx) + b_n \sin(nx)
    \end{equation}
    We are given a derivation for $a_0$ by starting with \cref{problem3:1} and taking the integral of both sides:
    \begin{equation}
        a_0 = \frac{1}{2\pi} \int_{-\pi}^{\pi}f(x) \, dx
    \end{equation}
    We use a similar method to derive the constants $a_m$ and $b_m$. Consider a fixed $m \geq 1$. We can now derive $a_m$ and $b_m$ from \cref{problem3:1}.
    To derive $a_m$, we start with
    \begin{equation}
        f(x) = a_0 + \sum_{n=1}^{\infty} a_n \cos(nx) + b_n \sin(nx) 
    \end{equation}
    Multiplying both sides by $\cos(mx)$ and integrating both sides yields
    \begin{equation}
        \int_{-\pi}^{\pi} f(x) \cos(mx) \, dx = \int_{-\pi}^{\pi} a_0 \cos(mx) \, dx + \int_{-\pi}^{\pi} \left( \cos(mx) \sum_{n=1}^{\infty} a_n \cos(nx) + b_n \sin(nx) \right) \, dx. 
    \end{equation}
    The integral with $a_0$ resolves to 0, so the derivation continues as
    \begin{align*}
        \int_{-\pi}^{\pi} f(x) \cos(mx) \, dx &= \int_{-\pi}^{\pi} \left( \cos(mx) \sum_{n=1}^{\infty} a_n \cos(nx) + b_n \sin(nx) \right) \, dx \\
        &= \int_{-\pi}^{\pi} \left( \sum_{n=1}^{\infty} a_n \cos(nx) \cos(mx) + b_n \sin(nx) \cos(mx) \right) \, dx \\
        &= \sum_{n=1}^{\infty} \left( \int_{-\pi}^{\pi} a_n \cos(nx) \cos(mx) + b_n \sin(nx) \cos(mx) \, dx \right) \\
        &= \sum_{n=1}^{\infty} \int_{-\pi}^{\pi} a_n \cos(nx) \cos(mx) + 0 \, dx
    \end{align*}
    Observing the integral, we see that $\cos(nx)\cos(mx)$ is 0 whenever $n \neq m$, and equals $\pi$ whenever $n=m$. Thus, the right hand side gets replaced with 
    \begin{equation}
        a_m \int_{-\pi}^{\pi} \cos^2(mx) \, dx = a_m \pi.
    \end{equation}
    Dividing both sides by $\pi$ yields the desired result:
    \begin{equation}
        a_m = \frac{1}{\pi} \int_{-\pi}^{\pi} f(x) \cos(mx) \, dx.
    \end{equation}

    We follow the same process to derive $b_n$, but instead of multiplying by $\cos(mx)$, we multiply by $\sin(mx)$ and integrate.
    \begin{align*}
        f(x) &= a_0 + \sum_{n=1}^{\infty} a_n \cos(nx) + b_n \sin(nx) \\
        \int_{-\pi}^{\pi} f(x) \sin(mx) \, dx &= \int_{-\pi}^{\pi} a_0 \sin(mx) \, dx + \int_{-\pi}^{\pi} \left( \sin(mx) \sum_{n=1}^{\infty} a_n \cos(nx) + b_n \sin(nx) \right) \, dx. \\
         &= \int_{-\pi}^{\pi} \left( \sum_{n=1}^{\infty} a_n \cos(nx) \sin(mx) + b_n \sin(nx) \sin(mx) \right) \, dx \\
         &= \sum_{n=1}^{\infty} \left( \int_{-\pi}^{\pi} a_n \cos(nx) \sin(mx) + b_n \sin(nx) \sin(mx) \, dx \right) \\
         &= \sum_{n=1}^{\infty} \int_{-\pi}^{\pi} b_n \sin(nx) \sin(mx) \, dx
    \end{align*}
    We see the same pattern where the terms of the sum equal 0 unless $n=m$. Since the integral resolves to $b_m\pi$, We land at
    \begin{equation}
        b_m = \frac{1}{\pi} \int_{-\pi}^{\pi} f(x) \sin(mx) \, dx.
    \end{equation}

    Thus, we have derived the following formulas for the constants $a_n$ and $b_n$:
    \begin{equation}
        a_m = \frac{1}{\pi} \int_{-\pi}^{\pi} f(x) \cos(mx) \, dx 
        \quad \text{and} \quad 
        b_m = \frac{1}{\pi} \int_{-\pi}^{\pi} f(x) \sin(mx) \, dx
    \end{equation}
    for all $m \geq 1$.
\end{problem}

\include{problem_4.tex}

\include{problem_5.tex}

\include{problem_6.tex}

\include{problem_7.tex}

\end{document}